\documentclass[10pt,a4paper]{article}
\usepackage{lmodern}
\usepackage[a4paper,top=18mm,bottom=17mm,left=18mm,right=18mm,headheight=21pt]{geometry}

\usepackage{microtype}
\usepackage{xcolor}
\usepackage{array,booktabs,tabularx,longtable}
\usepackage{enumitem}
\usepackage{titlesec}
\usepackage{fancyhdr}
\usepackage{hyperref}
\usepackage{lastpage}
\usepackage{amssymb}
\usepackage{tikz}

\definecolor{Navy}{HTML}{16324F}
\definecolor{Blue}{HTML}{286FA3}
\definecolor{Teal}{HTML}{23858C}
\definecolor{Pale}{HTML}{EAF2F7}
\definecolor{PaleTeal}{HTML}{E7F4F3}
\definecolor{Ink}{HTML}{25313B}
\definecolor{Muted}{HTML}{66727C}
\definecolor{Line}{HTML}{C9D5DE}
\definecolor{Warn}{HTML}{9A6110}

\hypersetup{colorlinks=true,linkcolor=Blue,urlcolor=Blue,pdftitle={V1 Robot Development and Validation Plan},pdfauthor={Jan Cuesta}}
\setlength{\parindent}{0pt}
\setlength{\parskip}{4pt}
\setlength{\emergencystretch}{2em}
\setlength{\tabcolsep}{5pt}
\renewcommand{\arraystretch}{1.23}
\setlist[itemize]{leftmargin=5mm,itemsep=1.5pt,topsep=2pt}

\titleformat{\section}{\Large\bfseries\color{Navy}}{}{0pt}{}
\titleformat{\subsection}{\large\bfseries\color{Blue}}{}{0pt}{}
\titlespacing*{\section}{0pt}{8pt}{5pt}
\titlespacing*{\subsection}{0pt}{7pt}{3pt}

\pagestyle{fancy}
\fancyhf{}
\lhead{\textcolor{Muted}{V1 Robot Development and Validation Plan}}
\rhead{\textcolor{Muted}{Jan Cuesta}}
\cfoot{\textcolor{Muted}{\thepage\ / \pageref{LastPage}}}
\renewcommand{\headrulewidth}{0.3pt}
\renewcommand{\headrule}{\hbox to\headwidth{\color{Line}\leaders\hrule height \headrulewidth\hfill}}

\newcommand{\cb}{\(\square\)}
\newcommand{\exitcheck}[1]{\vspace{2pt}\colorbox{PaleTeal}{\parbox{0.965\linewidth}{\textbf{Exit check:} #1}}}
\newcommand{\workstep}[2]{\subsection{\texorpdfstring{#1\quad #2}{#1\ #2}}}
\newcommand{\stagebanner}[2]{%
  \vspace{2pt}
  \noindent\colorbox{Navy}{\parbox{0.965\linewidth}{\color{white}\textbf{#1}\hfill #2}}
  \vspace{2pt}
}
\newcommand{\req}[1]{\texttt{#1}}

\begin{document}

\begin{titlepage}
\thispagestyle{empty}
\begin{tikzpicture}[remember picture,overlay]
  \fill[Navy] (current page.north west) rectangle ([yshift=-72mm]current page.north east);
  \fill[Teal] ([yshift=-72mm]current page.north west) rectangle ([yshift=-76mm]current page.north east);
\end{tikzpicture}

\vspace*{19mm}
{\color{white}\fontsize{27}{32}\selectfont\bfseries V1 ROBOT DEVELOPMENT\\[2mm]AND VALIDATION PLAN}\par
\vspace{3mm}
{\color{white}\Large Autonomous Semantic Navigation Robot}\par

\vspace{42mm}
{\Large\bfseries Jan Cuesta}\par
\vspace{2mm}
{\large Isaac Sim + ROS 2 Jazzy + Nav2}\par
\vspace{2mm}
{\color{Muted}Version 2.0 | Final planning baseline | 21 September 2026}\par

\vfill
\colorbox{Pale}{\parbox{0.91\textwidth}{
\textbf{Purpose}\par
Build and validate an indoor autonomous navigation system in simulation, then package it for controlled integration with the ASU robot. The robot must start reliably, load a saved map, localize, accept a destination, and reach it autonomously---or stop safely and report why it cannot.
}}

\vspace{8mm}
{\small\color{Muted}This plan follows an iterative systems-engineering workflow: needs and requirements, architecture, implementation, integration, verification, validation, and transition. Every work step advances only after its exit check passes.}
\end{titlepage}

\section{Plan control and completion boundary}

\textbf{V1 includes:} simulated differential-drive motion, a 2D LiDAR, wheel-derived odometry, mapping, saved-map localization, Nav2 navigation, a C++ mission manager, fault handling, logging, repeatable testing, and a documented interface for V2.

\textbf{V1 excludes:} RGB-D object recognition, semantic destination selection, manipulation, unbounded dynamic crowds, and proof of physical safety. Those require later project phases and real-hardware testing.

\begin{tabularx}{\textwidth}{>{\bfseries}p{31mm} X}
\toprule
Milestone & Required evidence \\
\midrule
Exploratory prototype & Isaac Sim, a lightweight mobile base, and basic movement are technically feasible. \\
Working demo & The robot maps a world, reloads the map, localizes, and reaches selected goals. \\
V1 complete & Behavior is repeatable; code, faults, startup, tests, evidence, and interfaces are documented. \\
V2 validated & Real sensors, motors, calibration, stopping behavior, timing, and physical navigation trials pass. \\
\bottomrule
\end{tabularx}

\textbf{Working rule:} an attractive demonstration is not a release. Advance when the exit check passes, retain failed results, and place code, configuration, maps, scene versions, and test evidence under version control.

\section{Development lifecycle}

\begin{tabularx}{\textwidth}{>{\bfseries\color{Navy}}p{29mm} X >{\raggedleft\arraybackslash}p{25mm}}
\toprule
Stage & Purpose & Work steps \\
\midrule
01 Definition & Establish intended use, requirements, risks, and hardware assumptions. & 1--4 \\
02 Architecture & Define interfaces, ownership, repository, and continuous quality controls. & 5--6 \\
03 Simulation & Build the world and integrate motion, sensors, frames, time, and odometry. & 7--10 \\
04 Autonomy & Create the map, localize, navigate, and execute missions. & 11--14 \\
05 Resilience & Handle faults, evaluate estimation, and automate evidence collection. & 15--17 \\
06 Validation & Freeze, test, correct, release, and prepare staged hardware integration. & 18--20 \\
\bottomrule
\end{tabularx}

\textbf{Iteration rule:} the stages describe dependency order, not a one-pass waterfall. Failed tests return the work to the relevant requirement, model, interface, or implementation step. Risks, requirements, testing, documentation, and configuration control continue throughout the project.

\newpage
\stagebanner{01 / DEFINITION AND GOVERNANCE}{Needs before construction}

\workstep{1.}{Write the Concept of Operations and freeze V1 scope}
\begin{itemize}
  \item[\cb] Define operator, intended indoor use, startup, shutdown, mapping, saved-map navigation, goal entry, cancellation, recovery, and manual-stop workflows.
  \item[\cb] Select one differential-drive simulation model, wheel odometry, one 2D LiDAR, and one lightweight test world. Record why the choices transfer to the intended ASU robot.
  \item[\cb] Define world geometry and difficulty: rooms, corridor, doorway widths, static obstacles, temporary blockage, start poses, and at least ten fixed goals.
  \item[\cb] Allow manual driving while mapping; require autonomous motion during navigation acceptance trials.
  \item[\cb] Record exclusions and assumptions, including floor, people, obstacle motion, network, lighting, and manual supervision.
\end{itemize}
\exitcheck{A concise ConOps lets another person explain how V1 is operated, what it must do, and what it intentionally does not do.}

\workstep{2.}{Baseline measurable requirements, traceability, and risks}
\begin{itemize}
  \item[\cb] Give every requirement an identifier and a measurable ``shall'' statement; classify it as functional, performance, interface, safety, or evidence.
  \item[\cb] Create a traceability table linking each requirement to its design component, verification method, test case, and final result.
  \item[\cb] Freeze definitions of success, failure, collision, timeout, manual intervention, acceptable pose error, and stopping response.
  \item[\cb] Establish timing and performance budgets: physics rate, real-time factor, topic rates, maximum data age, command latency, startup time, trial timeout, and sustained duration.
  \item[\cb] Create a lightweight risk register covering uncontrolled motion, stale commands, localization loss, stale sensors, incorrect geometry, simulation slowdown, thermal throttling, and version incompatibility.
\end{itemize}
\exitcheck{Every V1 requirement is testable and every major risk has an owner, mitigation, trigger, and planned test.}

\textbf{Initial requirement register}\par
\vspace{2pt}
\begin{tabularx}{\textwidth}{>{\ttfamily}p{18mm} X p{38mm}}
\toprule
ID & Requirement summary & Primary evidence \\
\midrule
NAV-01 & Reach valid goals autonomously on the saved map. & 30-trial campaign \\
LOC-01 & Localize after a clean restart from three selected starts. & Map-reuse tests \\
SAFE-01 & Stop predictably when velocity commands expire. & Command-loss test \\
SAFE-02 & Do not collide during the frozen acceptance set. & Collision log/video \\
PERF-01 & Maintain the frozen real-time-factor and topic-rate limits. & Performance log \\
ACC-01 & Meet frozen position and heading error tolerances. & Ground-truth comparison \\
DIAG-01 & Report clear mission success, cancellation, timeout, or failure. & Mission logs/tests \\
\bottomrule
\end{tabularx}

\newpage
\workstep{3.}{Confirm the physical robot interface}
\begin{itemize}
  \item[\cb] Obtain the exact ASU robot model, onboard computer, operating system, ROS distribution, access method, and responsible contact.
  \item[\cb] Record footprint, wheel radius/separation, maximum speeds, LiDAR/IMU models, sensor mounting frames, and available odometry.
  \item[\cb] Confirm velocity topic and message type, real drivers, command watchdog, emergency/manual stop, power-up behavior, and motor enable sequence.
  \item[\cb] Decide which processes run onboard and which run on the laptop; document expected network, clock, and QoS behavior.
  \item[\cb] Mark unknown values as provisional rather than silently treating simulation values as hardware facts.
\end{itemize}
\exitcheck{The intended simulation-to-hardware contract is known, and every unresolved mismatch is explicitly listed.}

\workstep{4.}{Establish a reproducible development environment}
\begin{itemize}
  \item[\cb] Record Ubuntu, NVIDIA driver, Isaac Sim, ROS 2 Jazzy, Nav2, SLAM Toolbox, RViz, compiler, CMake, and colcon versions.
  \item[\cb] Confirm that the selected Isaac Sim release and ROS 2 workflow are compatible; freeze working versions during the sprint.
  \item[\cb] Build and run a small C++ ROS 2 node; verify bidirectional ROS communication with Isaac Sim.
  \item[\cb] Document terminal sourcing, build, test, launch, shutdown, and clean-start commands.
  \item[\cb] Retain the 30-FPS safe launcher as the development baseline; treat viewport FPS separately from physics and real-time factor.
\end{itemize}
\exitcheck{A clean terminal can reproduce the build and communicate with the simulator using documented commands.}

\newpage
\stagebanner{02 / ARCHITECTURE AND PROJECT FOUNDATION}{Interfaces before integration}

\workstep{5.}{Define system architecture and interface control}
\begin{itemize}
  \item[\cb] Draw the node/component architecture for simulator, robot base, sensors, localization, Nav2, mission manager, logging, and safety timeout.
  \item[\cb] Create an interface table for topics, actions, services, message types, frame IDs, expected rates, QoS, time source, and publisher ownership.
  \item[\cb] Define the TF tree before sensor integration: \texttt{map -> odom -> base\_link -> sensors}, with documented platform-specific intermediate frames.
  \item[\cb] Assign exactly one publisher to each dynamic transform and one controlled path for motion commands.
  \item[\cb] Define which simulated components are replaced by hardware drivers in V2 without changing upper-level navigation interfaces.
\end{itemize}
\exitcheck{The ROS graph, TF ownership, command path, and simulation/hardware replacements are unambiguous before implementation.}

\workstep{6.}{Create the repository, packages, and continuous quality controls}
\begin{itemize}
  \item[\cb] Create the workspace under \texttt{~/projects}; initialize Git with README, license decision, issue list, change log, and tagged baselines.
  \item[\cb] Organize the description, simulation, bringup, navigation, mission, and test responsibilities as the six ROS packages defined in the architecture.
  \item[\cb] Store parameters in versioned configuration files; organize maps, worlds, launch files, mission files, results, and documentation.
  \item[\cb] Enable \texttt{ament\_lint}; add unit tests with the first production code and integration tests with the first cross-package interface.
  \item[\cb] Establish the repeatable quality command: build without warnings, run \texttt{colcon test}, inspect results, and perform a clean-start smoke test.
  \item[\cb] Exclude generated builds and large bags; document how large simulator assets are retrieved.
\end{itemize}
\exitcheck{The project rebuilds from source, package responsibilities are clear, and automated checks run from the beginning.}

\newpage
\stagebanner{03 / SIMULATION AND COMPONENT INTEGRATION}{Build the minimum valid system}

\workstep{7.}{Build and version the lightweight Isaac Sim acceptance world}
\begin{itemize}
  \item[\cb] Build the frozen floor, walls, corridor, doorways, localization features, and obstacle set using simple geometry and collision shapes.
  \item[\cb] Verify dimensions, robot clearance, friction, spawn poses, lighting, collision geometry, and obstacle mobility required by tests.
  \item[\cb] Use a mobile base approximating the physical platform; remove cameras, decorative assets, and unused sensor pipelines from V1.
  \item[\cb] Save the scene and dependencies; close Isaac Sim, reopen the stage, and verify the same configuration loads.
  \item[\cb] Establish a performance baseline for world plus robot before sensors and navigation are added.
\end{itemize}
\exitcheck{The complete scene reloads correctly, matches the frozen test geometry, and runs without prolonged stalls.}

\workstep{8.}{Connect motion to ROS and enforce command expiry}
\begin{itemize}
  \item[\cb] Connect ROS 2 velocity commands to the differential-drive controller using the interface contract.
  \item[\cb] Test forward, reverse, left rotation, and right rotation independently; verify axes, units, wheel order, radius, and separation.
  \item[\cb] Apply conservative velocity, acceleration, and deceleration limits appropriate to the acceptance-world clearance.
  \item[\cb] Maintain one controlled command path and prevent competing publishers from commanding the base unexpectedly.
  \item[\cb] Implement controller/base-level command expiry and verify predictable stopping after command loss.
\end{itemize}
\exitcheck{Motion matches commands and stale, absent, or terminated command sources produce the documented stop response.}

\workstep{9.}{Integrate sensors, simulation time, and transforms}
\begin{itemize}
  \item[\cb] Publish 2D LiDAR, wheel/joint measurements, odometry, simulation clock, and sensor transforms; add IMU data only if required for the V2 interface or planned comparison.
  \item[\cb] Validate timestamps, simulation-time use, message rates, maximum data age, QoS compatibility, and publisher ownership.
  \item[\cb] Verify in RViz that laser returns align with walls, sensor frames are correctly located, odometry direction is correct, and TF has no recurring errors.
  \item[\cb] Measure performance again after each sensor or bridge addition; record the change in real-time factor, CPU/GPU load, temperatures, and topic rates.
\end{itemize}
\exitcheck{The moving robot produces coherent, timely sensor data consistent with the architecture and performance budgets.}

\workstep{10.}{Validate transfer-oriented wheel odometry}
\begin{itemize}
  \item[\cb] Identify whether any Isaac example publishes perfect simulator pose; label it ground truth rather than odometry.
  \item[\cb] Use ground truth only for early debugging and independent error measurement; derive navigation odometry from simulated wheel motion.
  \item[\cb] Test straight lines, rotation, and a square trajectory; measure translation and heading errors against ground truth.
  \item[\cb] Repeat with controlled wheel-radius error, wheel-separation error, slip/noise, or bias appropriate to the risk plan.
  \item[\cb] Document the source, covariance assumptions, update rate, and transform ownership of every pose estimate.
\end{itemize}
\exitcheck{Navigation does not depend on perfect global simulator pose, and odometry performance is measured and documented.}

\newpage
\stagebanner{04 / MAPPING, LOCALIZATION, AND AUTONOMY}{Integrate only after components pass}

\workstep{11.}{Create and save the occupancy map}
\begin{itemize}
  \item[\cb] Configure SLAM Toolbox using the validated sensor, odometry, TF, and time interfaces.
  \item[\cb] Drive slowly through the world, complete useful loops, and inspect scan alignment continuously.
  \item[\cb] Investigate duplicated walls, distortion, discontinuities, and pose jumps before accepting the map.
  \item[\cb] Save the occupancy map and metadata; save the SLAM session separately only if mapping continuation is required.
  \item[\cb] Version the world, robot configuration, SLAM parameters, and map together.
\end{itemize}
\exitcheck{A usable map has been produced from simulated robot sensors and survives a clean restart.}

\workstep{12.}{Localize on the saved map after restart}
\begin{itemize}
  \item[\cb] Start saved-map navigation mode, load the map, and use AMCL for localization.
  \item[\cb] Supply an approximate initial pose and test at least three frozen starting positions and orientations.
  \item[\cb] Ensure that only the selected mapping or localization component publishes \texttt{map -> odom}.
  \item[\cb] Block mission acceptance until localization health meets the defined criterion.
  \item[\cb] Test an intentionally incorrect initial pose and document detection, recovery, or required operator action.
\end{itemize}
\exitcheck{The robot reliably localizes on the existing map without rebuilding it and reports when localization is not trustworthy.}

\workstep{13.}{Configure and tune autonomous navigation}
\begin{itemize}
  \item[\cb] Configure robot footprint, global and local costmaps, obstacle layers, inflation/clearance, planner, controller, and behavior plugins.
  \item[\cb] Set velocity/acceleration limits, goal and progress checkers, recovery behavior, lifecycle startup, and mission readiness checks.
  \item[\cb] Progress through open-space goals, turns, doorways, corridors, and multi-room routes.
  \item[\cb] Correct geometry, sensing, timing, and TF faults before compensating through planner/controller tuning.
  \item[\cb] Record each relevant tuning change and rerun affected routes to detect regressions.
\end{itemize}
\exitcheck{Representative routes complete repeatedly without manual steering, collision, or unexplained recovery loops.}

\workstep{14.}{Develop the C++ mission manager}
\begin{itemize}
  \item[\cb] Load named waypoints or a mission file, validate inputs, and wait for localization and Nav2 readiness.
  \item[\cb] Send goals through a ROS action client; track feedback, acceptance, rejection, completion, cancellation, and failure.
  \item[\cb] Implement explicit \texttt{idle}, \texttt{preparing}, \texttt{navigating}, \texttt{succeeded}, \texttt{failed}, and \texttt{cancelled} states.
  \item[\cb] Apply frozen timeouts, bounded retries, status reporting, and result logging.
  \item[\cb] Unit-test mission parsing and state transitions; integration-test rejection, cancellation, timeout, and retry limits.
\end{itemize}
\exitcheck{A predefined mission executes, cancels safely, and produces a clear, testable final result.}

\newpage
\stagebanner{05 / RESILIENCE, ESTIMATION, AND EVIDENCE}{Demonstrate failure behavior}

\workstep{15.}{Implement fault handling and safe stopping}
\begin{itemize}
  \item[\cb] Convert the risk register into explicit responses for stale LiDAR, odometry, TF, localization loss, blocked paths, unreachable goals, and overheating/performance limits.
  \item[\cb] Inject command loss, mission-process exit, Nav2 failure, simulator pause/reset, and selected sensor/transform failures.
  \item[\cb] Specify which conditions block startup, cancel a mission, stop motion, attempt bounded recovery, or require explicit restart.
  \item[\cb] Keep command-expiry protection at the base/controller level; do not depend solely on the mission process remaining alive.
  \item[\cb] Retain diagnostic evidence for each injected fault and verify the observed response against the requirement.
\end{itemize}
\exitcheck{Injected faults trigger the intended stop, recovery, or failure response without uncontrolled motion.}

\workstep{16.}{Evaluate sensor fusion only when justified}
\begin{itemize}
  \item[\cb] Confirm whether V2 requires locally configured wheel-odometry/IMU fusion or already provides an equivalent estimate.
  \item[\cb] If required, configure \texttt{robot\_localization} EKF; verify IMU axes, units, covariances, duplicated information, and transform ownership.
  \item[\cb] Compare wheel-only and fused estimates on identical trajectories using independent simulator ground truth.
  \item[\cb] Document improvements, regressions, parameter sensitivity, and the recommended V2 configuration.
  \item[\cb] Defer fusion if it adds no requirement-backed value; record the decision rather than adding complexity by default.
\end{itemize}
\exitcheck{Measured evidence supports either the proposed fusion configuration or a documented decision to defer it.}

\workstep{17.}{Complete automated system testing, logging, and metrics}
\begin{itemize}
  \item[\cb] Record scenario, start pose, goal, code/configuration versions, success, completion time, final pose error, and final mission state.
  \item[\cb] Log collisions, interventions, recoveries, sensor/TF faults, latency, topic rates, real-time factor, resource use, and thermal results.
  \item[\cb] Retain representative rosbag runs and important failures; separate true-pose error from estimated-pose metrics.
  \item[\cb] Automate end-to-end simulation scenarios while preserving unit and integration tests created throughout development.
  \item[\cb] Produce a requirement-to-test report showing that every frozen requirement has evidence or an explicit open disposition.
\end{itemize}
\exitcheck{Scenarios are repeatable, results are comparable across changes, and requirement coverage is visible.}

\newpage
\stagebanner{06 / ACCEPTANCE, RELEASE, AND HARDWARE TRANSITION}{Freeze before judging}

\workstep{18.}{Run the frozen acceptance campaign}
\begin{itemize}
  \item[\cb] Freeze code, world, map, parameters, tolerances, timeouts, sensor-age limits, and stopping responses before official testing.
  \item[\cb] Perform three clean cold starts and saved-map localization from three selected start poses.
  \item[\cb] Run ten fixed navigation goals three times, covering turns, doorways, corridors, and multi-room routes.
  \item[\cb] Test temporary blockage, unreachable goals, changed obstacles, imperfect odometry, and sensor/command loss.
  \item[\cb] Run a sustained session and retain all failures even when the aggregate success target passes.
\end{itemize}
\exitcheck{Every frozen requirement has a pass/fail result with traceable supporting evidence.}

\workstep{19.}{Investigate defects and perform regression testing}
\begin{itemize}
  \item[\cb] Prioritize uncontrolled motion, collisions, localization loss, crashes, startup failures, and timing violations.
  \item[\cb] Reproduce important failures; identify the violated requirement and plausible root cause.
  \item[\cb] Change one relevant factor at a time, document the rationale, and rerun the affected component/integration tests.
  \item[\cb] Rerun the final nominal campaign against the frozen release candidate.
  \item[\cb] Document at least two meaningful failure investigations and retain known limitations.
\end{itemize}
\exitcheck{No blocking defect remains, fixes have regression evidence, and accepted limitations are explicit.}

\workstep{20.}{Package the release and stage the V2 handoff}
\begin{itemize}
  \item[\cb] Provide separate simulation and hardware launch configurations with documented topic, message, frame, QoS, timing, and parameter contracts.
  \item[\cb] Remove simulation-only publishers from hardware launch; configure real-system time and clearly mark every unverified hardware value.
  \item[\cb] Include architecture, dependencies, startup/shutdown, maps, scene assets, mission procedures, calibration guidance, test evidence, demo video, and known limitations.
  \item[\cb] Tag the V1 code/configuration release and preserve the exact acceptance baseline.
  \item[\cb] Prepare the staged physical-integration checklist on the following page; do not proceed directly from simulation to autonomous driving.
\end{itemize}
\exitcheck{The V1 release is reproducible and ready for controlled physical integration, with remaining hardware verification clearly identified.}

\newpage
\section{Acceptance matrix}

The values below are project targets, not universal robotics standards. Confirm them during Step 2 against robot speed, sensor rates, available clearance, and ASU hardware constraints; then freeze them before Step 18.

\begin{longtable}{>{\bfseries}p{38mm}p{27mm}p{98mm}}
\toprule
Test & Requirement & Acceptance evidence \\
\midrule
\endhead
Cold startup & START-01 & Three clean starts without undocumented fixes. \\
Map reuse & LOC-01 & Saved map loads and localization succeeds from three chosen starts. \\
Nominal navigation & NAV-01 & Ten fixed goals repeated three times; initial target at least 27/30 successes without intervention. \\
Collision & SAFE-02 & Zero observed collisions in the acceptance set; any collision blocks release until investigated and dispositioned. \\
Goal accuracy & ACC-01 & Initial target: position error at most 0.20 m and heading error at most 10 degrees, measured against independent ground truth. \\
Temporary blockage & NAV-02 & Robot waits or replans and completes when a feasible route remains. \\
Unreachable goal & DIAG-01 & Robot stops and reports failure within the frozen timeout. \\
Command loss & SAFE-01 & Controller-level expiry produces the frozen stop response. \\
Sensor/TF loss & SAFE-03 & Detection, cancellation/stop behavior, and diagnostic output meet the frozen requirement. \\
Performance & PERF-01 & Real-time factor, topic rates, data age, latency, and thermal thresholds remain within the frozen budget. \\
Robustness & ROB-01 & Changed obstacles and imperfect odometry are tested; observed limitations are retained. \\
Sustained operation & PERF-02 & Thirty-minute session without unexplained crash, accumulating fault, or timing degradation. \\
\bottomrule
\end{longtable}

\section{Staged V2 physical-integration gate}

\begin{enumerate}[leftmargin=8mm,itemsep=2pt]
  \item[\cb] Confirm emergency stop, manual stop, power isolation, safe test area, supervision, and low-speed limits.
  \item[\cb] Verify real topics, timestamps, TF, LiDAR/IMU data, and network behavior with motors disabled.
  \item[\cb] Verify left/right motor direction and encoder sign with the platform safely lifted or otherwise restrained according to lab procedure.
  \item[\cb] Test command expiry and shutdown behavior before free motion.
  \item[\cb] Run very-low-speed forward, reverse, and rotation tests in an empty controlled area.
  \item[\cb] Calibrate wheel radius/separation, footprint, sensor extrinsics, and odometry/IMU behavior.
  \item[\cb] Validate stopping distance and clearance assumptions at the approved test speed.
  \item[\cb] Map a small controlled area, restart, and localize without autonomous motion.
  \item[\cb] Execute short open-area goals; progress later to corridors, doorways, and controlled obstacles.
  \item[\cb] Run at least 20 physical trials only after prior gates pass; initial target at least 90\% completion without intervention and zero uninvestigated collisions.
\end{enumerate}

\textbf{Important:} V1 simulation evidence does not validate physical braking, wheel slip, motor polarity, sensor calibration, networking, human interaction, or emergency-stop effectiveness.

\newpage
\section{Final release checklist}

\begin{tabularx}{\textwidth}{X X}
\cb\ Tagged code and pinned versions & \cb\ Reproducible build/test/launch instructions \\
\cb\ Baselined requirements and traceability & \cb\ Architecture and interface-control tables \\
\cb\ Saved map and complete scene dependencies & \cb\ Unit, integration, and system test results \\
\cb\ Acceptance trials and representative bags & \cb\ Two documented failure investigations \\
\cb\ Performance and thermal evidence & \cb\ Demo video and mission result examples \\
\cb\ Risk register and known limitations & \cb\ Staged V2 physical-integration checklist \\
\end{tabularx}

\section{Technical basis}

This project-scale plan adapts established engineering practice rather than claiming formal certification. The following sources support its lifecycle and technical ordering:

\begin{itemize}
  \item NASA, \href{https://www.nasa.gov/reference/systems-engineering-handbook/}{Systems Engineering Handbook}: stakeholder expectations, requirements, design, implementation, integration, verification, validation, transition, and continuous technical management.
  \item Nav2, \href{https://docs.nav2.org/rolling/configuration_and_development/first_time_robot_setup_guide/}{First-Time Robot Setup Guide}: transforms, robot description, odometry, perception sensors, mapping/localization, footprint, navigation plugins, and lifecycle bringup.
  \item ROS 2 Jazzy, \href{https://docs.ros.org/en/jazzy/The-ROS2-Project/Contributing/Developer-Guide.html}{Developer Guide} and \href{https://docs.ros.org/en/jazzy/The-ROS2-Project/Contributing/Quality-Guide.html}{Quality Guide}: documentation, unit/integration/system testing, linting, static analysis, and repeatable quality practices.
  \item NVIDIA Isaac Sim, \href{https://docs.isaacsim.omniverse.nvidia.com/latest/ros2_tutorials/robot_control/tutorial_ros2_navigation.html}{ROS 2 Navigation}: Isaac Sim, occupancy maps, Nav2, ROS actions, and programmatic goal execution.
\end{itemize}

\vfill
\colorbox{Pale}{\parbox{0.965\linewidth}{
\textbf{Project decision}\par
Isaac Sim replaces the older Gazebo assumption for V1. The simulator-specific bridge and model may change; the requirements, interfaces, tests, and navigation behavior remain the controlling project baseline.
}}

\end{document}
